% !TeX root = ../nc-quadric-classification.tex
% 2026-11-21 / v1.12
\section{可換 quadric の全一次方向を覆う二次 AS 被覆}
\label{sec:quadric-first-order-as}

$k$ は標数零、$A=k[\varepsilon]/(\varepsilon^2)$、
$P=k[x_0,x_1,x_2,x_3]$、$q=x_0x_3-x_1x_2$ とする。
$\epsilon_{ijkl}$ は $\epsilon_{0123}=1$ で定まる交代記号である。

\begin{proposition}[一次の Jacobian 被覆]
\label{prop:quadric-first-order-jacobian-cover}
$f\in P_2/kq$ を取り、$q_i=\partial_iq$、$f_i=\partial_if$ と書く。
$P$ 上の交代二重導分 $\gamma_f$ を
\[
 \gamma_f(x_i,x_j)=\sum_{k,l=0}^3\epsilon_{ijkl}q_kf_l
\]
で定める。$f$ の $kq$ による変更はこの式を変えない。
$\mathcal A_f=P\otimes_kA$ に積
\[
 a\star b=ab+\frac{\varepsilon}{2}\gamma_f(a,b)
\]
を入れると、これは次数別有限自由の二次 AS4 被覆となり、
ベクトル空間の同一視で指定した $q\in(\mathcal A_f)_2$ は中心正則である。
商 $\mathcal B_f=\mathcal A_f/(q)$ も全次数有限自由であり、
階数は $(n+1)^2$ である。
これらの平坦性・相対 Koszul 完全性・相対 Gorenstein 双対は底変更と両立する。
任意の Noetherian $A$-代数への底変更後にも左右 noetherian である。
\end{proposition}
\begin{proof}
証明は相対モジュライ別稿 \cite[\texttt{quadric_first_order} 節]{QuadricModuli2026} に置く。
ここでは後続の分類側から参照するため、仮定と主張を保持する。
\end{proof}

ここで $q$ は次数別加群 $P\otimes A$ 内の元である。
非可換の積記号を使って $x_0\star x_3-x_1\star x_2$ と書いた元とは、
一般に $\varepsilon$ の二次補正だけ異なる。

\begin{lemma}[全 bivector 方向の回復]
\label{lem:quadric-jacobian-bivectors}
$Q=\operatorname{Proj}(P/(q))$ の $x_0\ne0$ の chart で
$z=x_2/x_0$、$w=x_1/x_0$ と置く。
$\mathcal B_f$ の次数零局所化の一次交換子は
\[
 \{z,w\}=2\frac{f}{x_0^2}
        =2f(1,w,z,zw)
\]
である。従って $f\mapsto\gamma_f|_Q$ は
$P_2/kq\simeq H^0(Q,\bigwedge^2T_Q)$ という線形同型を与える。
\end{lemma}
\begin{proof}
証明は相対モジュライ別稿 \cite[\texttt{quadric_first_order} 節]{QuadricModuli2026} に置く。
ここでは後続の分類側から参照するため、仮定と主張を保持する。
\end{proof}

\begin{proposition}[Abel 変形の Hochschild 類と係数の規約]
\label{prop:quadric-jacobian-abelian-class}
$\operatorname{qgr}\mathcal B_f$ を閉 fiber $\operatorname{coh}Q$ の
pointed flat Abel 一次変形とみなしたとき、その Hochschild 類は、
各標準 affine chart 上で二コサイクル $\gamma_f/2$ により表される。
HKR の規約を
\[
 v\wedge w\longmapsto
 \left[(a,b)\longmapsto
   \frac{v(a)w(b)-w(a)v(b)}2\right]
\]
と固定すると、対応する bivector は
\[
 \pi_f=2f(1,w,z,zw)\,\partial_z\wedge\partial_w
 \quad (x_0\ne0)
\]
である。従って任意の指定された一次類
$u\leftrightarrow h(z,w)\partial_z\wedge\partial_w$
に対して、$f$ の制限を $h/2$ と取ればこの類を実現できる。
\end{proposition}
\begin{proof}
証明は相対モジュライ別稿 \cite[\texttt{quadric_first_order} 節]{QuadricModuli2026} に置く。
ここでは後続の分類側から参照するため、仮定と主張を保持する。
\end{proof}

\begin{lemma}[近恒等の補正と右 Zhang twist の規約]
\label{lem:quadric-first-order-root-twist}
$\delta\in\mathfrak{sl}_2\oplus\mathfrak{sl}_2$ を Segre 表示の $P_1$ に作用させ、
$\alpha=1+\varepsilon\delta$ と置く。
これは $\mathcal A_f$ と $\mathcal B_f$ の次数付き自己同型で、$\alpha(q)=q$ である。
$\mathcal C=\operatorname{qgr}\mathcal B_f$、$F=(1)$ とし、
その全軌道環が $\mathcal B_f$ に一致すると仮定する。
右加群に作用を $\alpha^{-1}$ で twist する自己同値を $G_\alpha$ とすれば、
標識 $G_\alpha(\pi\mathcal B_f)\simeq\pi\mathcal B_f$ を
$b\mapsto\alpha(b)$ で取ることにより、
$G_\alpha$ の全軌道環上の作用は $\alpha$ となる。
従って $s=G_\alpha F$ の右規約の全軌道環は
\[
 a*_{s}b=\alpha^{\deg b}(a)*_Fb
\]
という右 Zhang twist $\mathcal B_f^\alpha$ であり、
$\mathcal A_f^\alpha$ はその相対 AS4 被覆を与える。
\end{lemma}
\begin{proof}
証明は相対モジュライ別稿 \cite[\texttt{quadric_first_order} 節]{QuadricModuli2026} に置く。
ここでは後続の分類側から参照するため、仮定と主張を保持する。
\end{proof}


\subsection{可換 quadric の全一次変形に対する適合表示}
\label{subsec:quadric-first-order-bridge}

$k$ は標数零、$A=k[\varepsilon]/(\varepsilon^2)$ とする。
$Q=\mathbf P^1_k\times\mathbf P^1_k$、$L=\mathcal O_Q(1,1)$ と置き、
命題~\ref{prop:quadric-first-order-jacobian-cover} の
$\mathcal B_f=\mathcal A_f/(q)$ を用いる。
以下の最初の補題では、変形された圏が $\mathfrak H_2$ に属するとは仮定しない。

\begin{lemma}[商圏の平坦性と全次数の飽和]
\label{lem:quadric-first-order-flat-saturated}
$\mathcal G_A=\operatorname{QGr}\mathcal B_f$、
$\mathcal C_A=\operatorname{qgr}\mathcal B_f$、
$\mathcal O_A=\pi\mathcal B_f$、$F=(1)$ と置く。
このとき $(\mathcal C_A,\mathcal O_A)$ は
$(\operatorname{coh}Q,\mathcal O_Q)$ の flat pointed Abel 変形であり、
$\mathcal O_A$ は $A$ 上平坦である。
自然な射は全次数で同型
\[
 \mathcal B_f\xrightarrow{\ \sim\ }
 \bigoplus_{n\ge0}\operatorname{Hom}_{\mathcal C_A}
       (\mathcal O_A,F^n\mathcal O_A)
\]
となる。また $\mathcal O_A$ は相対例外的である。
\end{lemma}
\begin{proof}
証明は相対モジュライ別稿 \cite[\texttt{quadric_first_order} 節]{QuadricModuli2026} に置く。
ここでは後続の分類側から参照するため、仮定と主張を保持する。
\end{proof}

\begin{lemma}[平方根補正の全近恒等方向は次数付き自己同型から来る]
\label{lem:quadric-first-order-correction-geometric}
$\mathcal G_A$ の finite foundation 表示を取り、
その近恒等な derived Picard 群の核を $K_A$ と書く。
命題~\ref{prop:artin-serre-square-correction} の補正元は、
ある $\delta\in\mathfrak{sl}_2\oplus\mathfrak{sl}_2$ と
$\alpha=1+\varepsilon\delta$ に対する $G_\alpha$ で表される。
ここで $G_\alpha$ は右加群の作用を $\alpha^{-1}$ で twist する関手である。
\end{lemma}
\begin{proof}
証明は相対モジュライ別稿 \cite[\texttt{quadric_first_order} 節]{QuadricModuli2026} に置く。
ここでは後続の分類側から参照するため、仮定と主張を保持する。
\end{proof}

\begin{theorem}[可換 quadric の全一次方向は相対 $\mathfrak H_2$ に入る]
\label{thm:quadric-all-first-order-h2}
$u\in\operatorname{HH}^2(Q)$ を任意に取る。
この一次類を持つ $(\operatorname{coh}Q,\mathcal O_Q)$ の flat pointed Abel 変形には、
$s_0=(-)\otimes\mathcal O_Q(1,1)$ の持上げ $s_A$ と相対 AS4 被覆が存在し、
定義~\ref{def:relative-h2} の全有限型基底変更を含む相対族を構成できる。
従って、この点での $\mathfrak H_2$ から pointed Abel 変形への一次忘却写像は全射である。
\end{theorem}
\begin{proof}
証明は相対モジュライ別稿 \cite[\texttt{quadric_first_order} 節]{QuadricModuli2026} に置く。
ここでは後続の分類側から参照するため、仮定と主張を保持する。
\end{proof}


\begin{corollary}[可換点の内在的な接空間]
\label{cor:quadric-intrinsic-tangent}
可換対象 $x=(\operatorname{coh}Q,\mathcal O_Q)$ において、
一次変形の忘却により自然な同定
\[
 T_x\mathfrak H_2\simeq\operatorname{HH}^2(Q)
      \simeq H^0(Q,\mathcal O_Q(2,2))
\]
を得る。従って内在的接空間は九次元である。
反 Serre 関係式を表す点を $\bar x\in[X_3/\operatorname{PGL}_9]$ と書けば、
\[
 \dim H^0T_{\bar x}[X_3/\operatorname{PGL}_9]=15,
 \qquad \dim\ker(T_{\bar x}\to\operatorname{HH}^2(Q))=6.
\]
後者の六次元は canonical 反 Serre 偏極からのずれであり、
$\mathfrak H_2$ の追加変形方向ではない。
\end{corollary}
\begin{proof}
証明は相対モジュライ別稿 \cite[\texttt{quadric_first_order} 節]{QuadricModuli2026} に置く。
ここでは後続の分類側から参照するため、仮定と主張を保持する。
\end{proof}

全 Artin 変形の認識は系~\ref{cor:quadric-formal-deformation-equivalence} で得られる。
これらの形式的結果だけから内在的 $\mathfrak H_2$ 全体の代数性は従わない。
